\chapter{Lateral boundaries and nesting}\label{ch:bdy}

A limited-area model needs information from outside its domain. WRF's outer domain takes it from the boundary file
of \cref{ch:icbc}; each nest takes it from its parent during the run. This chapter describes:
\begin{itemize}
  \item the specified and relaxation zones of the outer domain;
  \item the other lateral boundary types;
  \item how a nest is forced by its parent, and how it can feed back.
\end{itemize}

\section{Specified and relaxation zones}

The outer $w_b$ rows and columns of the domain (\code{spec_bdy_width}, 5 by default) form the boundary zone,
divided into two parts:
\begin{itemize}
  \item the \emph{specified zone}, the outermost \code{spec_zone} points (1 by default);
  \item the \emph{relaxation zone}, the next \code{relax_zone - 1} points (\code{relax_zone} = 4 by default, so three
    points).
\end{itemize}

\paragraph{Specified zone.} There the prognostic variables follow the boundary data exactly. The slow tendency is
replaced by the time derivative of the boundary values
(\src{share/module_bc.F}{spec_bdytend}). In every acoustic sub-step the zone is updated with that derivative
(\src{share/module_bc.F}{spec_bdyupdate}), so it is never computed from the model's own equations.

\paragraph{Relaxation zone.} There the model is nudged toward the boundary data. With $\psi_{LS}$ the boundary value
at the current time, $\psi_{LS} = \psi_b(t_n) + (t - t_n)\,\partial_t\psi_b$ (the code's
\code{field_bdy + dtbc*field_bdy_tend}), and $n$ the distance in grid points from the edge, the tendency gets
\begin{equation}
  \left.\pd{\psi}{t}\right|_n \mathrel{+}= F_1(n)\,(\psi_{LS} - \psi) - F_2(n)\,\Delta^2(\psi_{LS} - \psi),
  \label{eq:relax}
\end{equation}
where $\Delta^2$ is the five-point Laplacian
$\Delta^2 f = f_{i-1,j} + f_{i+1,j} + f_{i,j-1} + f_{i,j+1} - 4 f_{i,j}$. The weights decrease linearly from the
edge inward, as set in \code{dyn_em/module_bc_em.F}:
\begin{equation}
  F_1(n) = \frac{0.1}{\Delta t}\,\frac{n_s + n_r - n}{n_r - 1}, \qquad
  F_2(n) = \frac{1}{50\,\Delta t}\,\frac{n_s + n_r - n}{n_r - 1},
\end{equation}
with $n_s$ = \code{spec_zone} and $n_r$ = \code{relax_zone}. An optional exponential decay (\code{spec_exp})
multiplies both. The relaxation is applied to the mass-coupled variables: \code{mass_weight} couples them first, and
\src{share/module_bc.F}{relax_bdytend_core} evaluates \cref{eq:relax} strip by strip (south, north, west, east).

\begin{algorithm}[ht]
\caption{The southern strip of \code{relax_bdytend_core} for one field.}
\label{alg:relax}
\begin{algorithmic}[1]
\For{$j$ from the inner edge of the specified zone to the inner edge of the relaxation zone}
  \State $n \gets j - j_{\mathrm{bs}}$ \Comment{distance from the edge}
  \For{$k$, then $i$ inside the strip}
    \State $f_0 \gets \psi_{LS}(i,j) - \psi(i,j)$; \quad $f_1, f_2, f_3, f_4$ the same at $i\mp1$, $j\mp1$
    \State $T(i,k,j) \gets T(i,k,j) + F_1(n)\,f_0 - F_2(n)\,(f_1 + f_2 + f_3 + f_4 - 4 f_0)$
  \EndFor
\EndFor
\end{algorithmic}
\end{algorithm}

\begin{portnote}
The four strips of \cref{alg:relax} overlap at the corners, and each strip runs over its own index ranges. The GPU
version keeps one kernel per strip, in the original order (Template G, kernels K-RLX-YS, -YE, -XS, -XE). Merging the
strips into one launch is a later, bit-neutral optimization (\cref{ch:fusion}).
\end{portnote}

\section{Other boundary types}

\begin{description}
  \item[Open] (\code{open_xs}, \dots) Radiation conditions let waves leave the domain: normal momentum is extrapolated
    from inside and advected out, and inflow uses zero-gradient values. The idealized fire case (S-3M) uses them.
  \item[Periodic] (\code{periodic_x}, \code{periodic_y}) The halo of one side is filled from the other.
  \item[Symmetric] (\code{symmetric_xs}, \dots) Mirror conditions.
  \item[Nested] As specified, with values from the parent.
\end{description}
Physical boundary conditions in the code are applied by
\src{share/module_bc.F}{set_physical_bc3d} and \code{set_physical_bc2d}. On a single rank they fill the halo cells of
the domain edges. Their $x$ part must run before their $y$ part, because the $y$ copies read the corners that the $x$
copies wrote.

\section{Nesting}

A nest is a finer grid inside a parent. It runs \code{parent_time_step_ratio} steps (9 in the case) for each parent
step. The recursive driver \src{frame/module_integrate.F}{integrate} does, for each parent step:
\begin{enumerate}
  \item one step of the parent (\code{solve_interface});
  \item for each nest:
    \begin{itemize}
      \item force its boundaries from the parent (\code{med_nest_force});
      \item integrate it over the parent's step (a recursive call of \code{integrate});
      \item with two-way nesting, feed its solution back into the parent (\code{med_nest_feedback}).
    \end{itemize}
\end{enumerate}
The case runs one-way (\code{feedback = 0}), so the feedback step is skipped.

\subsection{Forcing a nest}

\src{share/mediation_force_domain.F}{med_force_domain} gives a nest its boundary values for the coming parent step.
\begin{algorithm}[ht]
\caption{Forcing one nest for one parent step (\code{med_force_domain}).}
\label{alg:force}
\begin{algorithmic}[1]
\State allocate an \emph{intermediate} domain: the part of the parent that covers the nest, on the parent's grid
\State couple the parent's and the nest's prognostic fields with their column mass and map factors (\code{couple_or_uncouple_em}, couple)
\State copy the parent's fields that the nest needs into the intermediate domain (\code{interp_domain_em_part1}); on several ranks this is a redistribution
\State interpolate from the intermediate domain to the nest's boundary zone; compute the nest's boundary values and their tendencies over the parent step (\code{force_domain_em_part2}, \code{bdy_interp})
\State uncouple both domains again
\State free the intermediate domain
\end{algorithmic}
\end{algorithm}

The interpolation is selected by \code{interp_method_type}: bilinear, or the default \code{sint}, a smoothed
interpolation that conserves averages. The fields to interpolate, and how, are declared in the Registry with the
\code{f} (force) and \code{d} (down) flags (\cref{ch:arch}). Among them is \code{o3rad}, the ozone climatology of the
radiation, which the nest takes in full from the parent.

\paragraph{Couple, then uncouple.} Coupling multiplies a field by $\mu_d/m$; uncoupling multiplies by the stored
reciprocal. In floating point $x\cdot a\cdot(1/a) \ne x$ in general, so a couple--uncouple pair changes the last bits
of the nest's and parent's state at every parent step. The CPU reference does this, so the GPU must do it too, in the
same order, on the device (\cref{ch:subsys}).

\subsection{Feedback and smoothing}

With \code{feedback = 1}, after the nest has caught up, its solution is averaged onto the parent points it covers
(\code{med_nest_feedback}). The result is optionally smoothed (\code{smooth_option}). This makes the parent see the
nest's finer structures; it also makes the two domains depend on each other at every parent step.

\begin{portnote}
For the GPU, nest forcing is a data-movement problem more than a compute problem:
\begin{itemize}
  \item the interpolation runs on the host;
  \item the parent's fields that the host interpolation reads must be copied from the device;
  \item the nest's boundary arrays and \code{o3rad}, which it writes, must be copied back.
\end{itemize}
The port first copies both domains' whole state around the call (a correct but slow bridge). It then replaces this
with exact lists of fields and rows: a j-slab of the parent and strips of the nest (\cref{ch:subsys}).
\end{portnote}
